\subsection{Data Control Language}
\begin{frame}[fragile]{GRANT and REVOKE (DCL)}
  \begin{block}{Controlling Access to Database Objects}
\begin{lstlisting}[style=sqlstyle]
-- Grant all privileges on a database to a user
GRANT ALL ON supermarket.* TO 'admin'@'localhost';

-- Grant only SELECT on the whole database
GRANT SELECT ON supermarket.* TO 'reader'@'%';

-- Grant specific privileges on a single table
GRANT INSERT, UPDATE ON supermarket.product
    TO 'editor'@'localhost';

-- Revoke a privilege
REVOKE INSERT ON supermarket.product
    FROM 'editor'@'localhost';
\end{lstlisting}
  \end{block}

  \begin{block}{Principle of Least Privilege}
    \begin{itemize}
      \item Grant only the \textbf{minimum permissions} a user or
            application actually needs to do its job.
      \item \textbf{Never} run application code as the \texttt{root} (or
            \texttt{sa}) database user.
      \item \texttt{'user'@'\%'} allows connections from \textbf{any host};
            \texttt{'user'@'localhost'} restricts to the same machine only.
            Prefer the more restrictive option wherever possible.
    \end{itemize}
  \end{block}
\end{frame}
